\subsection{Principal findings}
Logistic regression achieved the strongest pooled discrimination in the historical fixed and fixed-feature tuned comparisons. Within the tuned families, it also achieved the highest F1 and matched XGBoost's sensitivity. The historical fixed random forest had higher sensitivity than its tuned counterpart, while tuning improved XGBoost's descriptive AUC and classification outcomes. LightGBM and stacking offered the highest tuned specificity. These trade-offs show that an inner-AUC search objective does not necessarily improve every outer metric, and that added complexity confers no uniform advantage. Dependency changes prevent attributing historical before/after differences solely to tuning; the observed rankings do not establish statistically significant superiority.

The updated base-tuned stack subsequently achieved AUC 0.8762 and F1 0.8133, with 26 false negatives, improving descriptively over its historical tuned counterpart. Jointly changing the base learners, their tuning, feature-count search and sampler scheduling prevents isolating the cause. The standalone candidates have not completed the same expanded feature-count search; shared outer exclusions and a common inner-CV budget do not by themselves make those configurations fully controlled comparators. Frequent selection of all 38 predictors also argues against presuming that feature reduction necessarily helps this architecture.

The findings accord with literature showing that optimisation and ensemble complexity do not consistently improve performance on Gallstone-1. Published scores cannot be compared directly with these LOPO estimates because training sizes, feature-selection procedures, thresholds, and assessment samples differ. Likewise, errors recurring across model families suggest shared limitations in the available measurements or fitted relationships, but retrospective inspection cannot determine a record's clinical circumstances or establish incorrect labels.

The additional SHAP and LIME analysis offers a visual account of model reliance and confident errors. Its global summary uses only 20 excluded records, and local examples were selected after observing errors. Modest LIME fidelity limits the strength of patient-specific interpretation. Neither method validates biological mechanisms or resolves errors; larger explanatory samples and stability analyses would be needed to assess robustness.

\subsection{Limitations}
The analysis relies on a small cohort from one clinical setting. Its near-balanced case distribution does not establish predictive values or assessment workload in populations with different prevalence. Availability of detailed bioimpedance measurements may further constrain transferability. Source-row identifiers are not verified person identifiers; the LOPO interpretation assumes one record per person, and repeated-person data would require grouped exclusion.

Fold-specific preprocessing prevents direct use of excluded records during fitting, but the analysis evaluates fixed pipelines and within-family nested tuning rather than the entire model-development process. The dataset and modelling choices had been explored previously; consequently, LOPO is an internal comparison rather than an untouched confirmatory assessment. Retrospective selection of the strongest family can introduce optimism, and the heavily overlapping training folds are not independent experiments. The single fixed-seed run does not quantify random-state variability or furnish confidence intervals for model differences.

The study did not establish calibrated probabilities, optimise a clinical operating threshold, or demonstrate external generalisation or clinical benefit. Feature selection may also vary across folds because correlated measurements can substitute for one another. Aggregate metrics and recurring error examples therefore require cautious interpretation and do not establish causal relationships.

\subsection{Implications for further research}
The implemented tuning nests parameter selection within outer training sets. Further work should extend this boundary to model-family choice, feature-count decisions, calibration, and operating-threshold selection. A sensitivity-oriented threshold should be assessed jointly with specificity and assessment workload. Grouped feature ablation and prespecified subgroup evaluation could investigate measurement burden and uneven errors. An untouched external cohort is required to assess transfer beyond this dataset. Prospective disease-onset prediction would require a different study design and outcome definition.
